\section[4. Recursion - {\it 재귀}]{4. Recursion}

\subsection{Recursion}

\textbf{재귀}\textit{\textsuperscript{Recursion}}은 함수가 자기 자신을 다시 호출하여 문제를 해결하는 프로그래밍의 기법이다.\\
반복문과 같이 반복적으로 작업을 수행하나 동일한 작업 수행하는 코드를 보다 간결하게 작성할 수 있다. 일부 문제를 해결하기 위해 더 작은 문제로 분할한 뒤에, 이들을 해결하여 전체 문제를 해결하는 과정을 거친다.\\
함수가 자신을 호출하면서 호출 횟수를 제한하는 조건문이 필수적이며 이러한 조건문이 없을 경우, 무한한 재귀 호출이 발생할 가능성이 있기에 주의해야 한다.
\begin{itemize}
\item 장점: 일부 문제에서 반복문에 비해 코드가 간결해진다. 그래프 탐색과 같은 여러 문제들의 해결에 용이하다.
\item 단점: 자기 자신 호출하는 과정에서 스택 메모리Stack Memory를 사용하기에 메모리 사용량이 증가하며 반복문보다 함수 호출과 반환에 대한 오버헤드(함수 호출시에 발생하는 추가적인 작업량)가 크기 때문에 실행 속도가 저하된다.
\end{itemize}

하노이 탑 문제는 3개의 기둥에 크기가 다른 원판들이 여러 개 있을 때, 한번에 하나의 원판을 옮겨서 첫 번째 기둥에 있는 모든 원판을 세 번째 기둥으로 옮기는데 걸리는 최소 반복 횟수를 구하는 문제이다. 이때, 작은 원판 위에 큰 원판을 옮길 수 없으며, 원판을 옮기는 전체 과정을 출력할 수 있어야 한다.
\begin{itemize}
\item 재귀적 해법: 초기 상태로 1번 기둥에 n개의 크기가 다른 원판이 놓여져 있다고 하자. 상위 n-1개 원판을 2번째 기둥으로 옮긴 다음에, n번째 원판을 3번째 기둥으로 옮기고, 상위 n-1개 원판을 3번째로 옮기면 된다.
\end{itemize}
